\documentclass[a4paper,11pt]{article}
\usepackage[a4paper,left=25mm,right=25mm,top=25mm,bottom=25mm]{geometry}
\usepackage[T1]{fontenc}
\usepackage[utf8]{inputenc}
\usepackage{lmodern}
\usepackage{setspace}
\usepackage{amsmath,amssymb,mathtools}
\usepackage{hyperref}
\usepackage{enumitem}
\usepackage{longtable}
\usepackage{booktabs}
\usepackage{array}
\usepackage{float}

\setlength{\parskip}{0.5em}
\setlength{\parindent}{0pt}
\setlength{\emergencystretch}{3em}
\setstretch{1.2}
\renewcommand{\arraystretch}{1.15}

\title{Comprehensive Technical Documentation of the Local Transformer Based Small Language Model Project}
\author{Final Year Project Documentation}
\date{\today}

\begin{document}
\sloppy
\maketitle

\tableofcontents
\newpage

\section{Introduction and System Intent}
This project is an end-to-end local Question Answering (QA) pipeline built on top of a small Transformer language model and a retrieval fallback mechanism.
The intent is to convert arbitrary domain PDFs into a queryable knowledge system without relying on external hosted LLM APIs.

The system is engineered around the following practical principle:
\begin{quote}
``Use generation when it is good enough, and use retrieval when generation is unreliable.''
\end{quote}

This dual strategy is important because local small models trained from scratch on limited corpora can be informative but not consistently fluent.
Retrieval improves factual reliability while generation preserves flexibility.

\subsection{High-Level Pipeline}
\[
\begin{aligned}
\text{PDFs} &\rightarrow \text{Raw Extracted Text} \rightarrow \text{Clean Corpus} \rightarrow \text{Sentence Bank} \\
&\rightarrow \text{Synthetic QA Pairs} \rightarrow \text{Tokenizer} \rightarrow \text{Transformer Training} \\
&\rightarrow \text{Inference + Retrieval Fallback}
\end{aligned}
\]

\subsubsection{Layman explanation of the pipeline arrow chain}
Think of the pipeline like a factory:
\begin{enumerate}[leftmargin=1.5em]
    \item \textbf{PDFs}: raw textbooks/notes (unstructured).
    \item \textbf{Raw Extracted Text}: machine-readable text pulled from PDFs.
    \item \textbf{Clean Corpus}: corrected/normalized text.
    \item \textbf{Sentence Bank}: clean text split into individual statements.
    \item \textbf{Synthetic QA Pairs}: sentence statements reshaped as question-answer examples.
    \item \textbf{Tokenizer}: text is converted into token IDs.
    \item \textbf{Transformer Training}: model learns patterns in token sequences.
    \item \textbf{Inference + Retrieval}: answer is produced by model; fallback retrieval is used if quality is weak.
\end{enumerate}

\section{Design Objectives and Constraints}
\subsection{Objectives}
\begin{enumerate}[leftmargin=1.5em]
    \item Build an offline, local QA stack with reproducible scripts.
    \item Support arbitrary PDF subjects, not a single fixed domain.
    \item Keep training and inference transparent and debuggable.
    \item Ensure data and model lifecycle can be repeated after adding new PDFs.
\end{enumerate}

\subsection{Constraints}
\begin{itemize}[leftmargin=1.5em]
    \item Source PDFs can contain OCR and formatting artifacts.
    \item The model is small and data is limited compared to industrial pretraining.
    \item Tokenization and text cleaning quality strongly affect final answers.
    \item Local compute limits model depth/width and training time.
\end{itemize}

\section{Repository Modules and Their Roles}
\begin{longtable}{>{\raggedright\arraybackslash}p{0.27\linewidth} >{\raggedright\arraybackslash}p{0.68\linewidth}}
\toprule
\textbf{Module/File} & \textbf{Detailed responsibility} \\
\midrule
\texttt{textextraction.py} & Scans \texttt{data/raw/} for PDFs; extracts each to text with \texttt{pdfminer.six}; writes outputs to \texttt{data/extracted/}. \\
\texttt{clean\_text.py} & Applies corpus normalization rules to all extracted text files and writes merged output to \texttt{data/cleaned/corpus\_clean.txt}. \\
\texttt{sentence\_split.py} & Splits merged corpus into sentence-level lines using NLTK; includes fallback behavior if tokenizer resources are unavailable. \\
\texttt{create\_dataset.py} & Converts sentence lines into synthetic \texttt{Ask/Answer} records and generates multiple question variants (definition prompts plus short explain/describe alternatives) per sentence. \\
\texttt{train\_tokenizer.py} & Trains SentencePiece tokenizer over the training text and exports tokenizer model/vocabulary files. \\
\texttt{model\_def.py} & Defines Transformer architecture: embeddings, positional encoding via learned embeddings, masked self-attention stack, and output projection. \\
\texttt{train.py} & Loads tokenizer and dataset, inserts EOS between records, applies a 90/10 train-validation split, trains with warmup+cosine LR schedule, gradient clipping, and early stopping, then saves the best checkpoint to \texttt{model.pth}. \\
\texttt{generate.py} & Runs autoregressive generation, weak-answer checks, hybrid retrieval scoring, question-type routing (factual vs definition/open-ended), answer trimming, and grammar cleanup. \\
\texttt{run\_pipeline.py} & Orchestrates the full pipeline with local-venv Python discovery and supports \texttt{--epochs}, \texttt{--skip-train}, and \texttt{--skip-tokenizer}. \\
\bottomrule
\end{longtable}

\section{Data Pipeline in Full Detail}
\subsection{Stage 1: PDF Text Extraction}
\subsubsection{Input/Output Contract}
\begin{itemize}[leftmargin=1.5em]
    \item Input directory: \texttt{data/raw/}
    \item Input type: \texttt{*.pdf}
    \item Output directory: \texttt{data/extracted/}
    \item Output type: one \texttt{.txt} file per PDF
\end{itemize}

\subsubsection{Operational Mechanics}
For each PDF file $D_i$, extraction computes:
\[
T_i = \texttt{extract\_text}(D_i)
\]
and writes $T_i$ to a path in \texttt{data/extracted/}.

Before writing new outputs, existing \texttt{.txt} files in \texttt{data/extracted/} are cleared so each run reflects only the current PDF set.

\subsubsection{Potential Noise Sources}
\begin{itemize}[leftmargin=1.5em]
    \item Header/footer duplication
    \item Broken ligatures and Unicode symbols
    \item Irregular line breaks from multi-column layouts
    \item Residual page artifacts and references
\end{itemize}

\subsection{Stage 2: Corpus Cleaning}
\subsubsection{Normalization Rules}
Given raw extracted text $R$, cleaning function $C(\cdot)$ applies regex transforms:
\begin{align*}
R_1 &= \text{collapse repeated newlines} \\
R_2 &= \text{remove page-number patterns} \\
R_3 &= \text{collapse whitespace runs}
\end{align*}
Then:
\[
C(R) = \text{strip}(R_3)
\]
All cleaned documents are concatenated into:
\[
\mathcal{C} = C(R_1) \oplus C(R_2) \oplus \cdots \oplus C(R_n)
\]
where $\oplus$ denotes newline-join composition.

\subsubsection{Term-by-term explanation and intuition}
Meaninng of each line:
\begin{itemize}[leftmargin=1.5em]
    \item $R$: original messy text.
    \item $R_1$: after fixing repeated line breaks.
    \item $R_2$: after removing page labels like ``Page 14''.
    \item $R_3$: after collapsing random extra spaces.
    \item $C(R)$: final cleaned result after trimming edges.
\end{itemize}

Why write this in steps? Because each transform has one clear purpose. In debugging, if output looks wrong, we can inspect intermediate states ($R_1, R_2, R_3$) instead of guessing.

Simple derivation logic:
\begin{enumerate}[leftmargin=1.5em]
    \item Apply transform A to reduce one kind of noise.
    \item Apply transform B to reduce another kind.
    \item Repeat until text is stable enough for tokenization.
\end{enumerate}

\subsubsection{Why this stage matters mathematically}
Tokenization and sentence splitting are sensitive to symbol and whitespace entropy.
Cleaning reduces noisy token variants, effectively shrinking unnecessary vocabulary branches and improving token frequency concentration.

\subsection{Stage 3: Sentence Segmentation}
The cleaned corpus is segmented into sentence candidates:
\[
\mathcal{S} = \{s_1, s_2, \ldots, s_m\}
\]
Primary splitter uses NLTK's statistical sentence boundary detector.
Fallback splitter uses punctuation regex when NLTK resources are missing.

\subsection{Stage 4: Synthetic QA Dataset Construction}
\subsubsection{Record Template}
The generated record format is:
\begin{verbatim}
Ask: <question>
Answer: <source sentence>
\end{verbatim}

\subsubsection{Question Heuristic}
For each sentence $s$:
\begin{itemize}[leftmargin=1.5em]
    \item If sentence matches definition-style pattern, construct \texttt{What is ...?}/\texttt{What are ...?}
    \item Add variants such as \texttt{Define ...}, \texttt{Explain ...}, and \texttt{What do you mean by ...?}
    \item Else construct fallback prompts such as \texttt{Explain: ...?} and \texttt{Describe this: ...}
    \item Remove duplicate question variants while preserving order
\end{itemize}

\subsubsection{Rationale}
This converts unstructured sentence corpus into instruction-like text pairs that align with prompt format used during inference.

\section{Tokenizer Theory and Implementation}
\subsection{SentencePiece Role}
Tokenizer maps text to discrete IDs:
\[
\tau: \text{string} \rightarrow \{0,1,\ldots,V-1\}^*
\]
and decoder maps back:
\[
\tau^{-1}: \{0,1,\ldots,V-1\}^* \rightarrow \text{string}
\]

\subsubsection{What each symbol means}
\begin{itemize}[leftmargin=1.5em]
    \item $\tau$: ``text to token IDs'' function.
    \item $\tau^{-1}$: ``token IDs back to text'' function.
    \item $V$: tokenizer vocabulary size.
    \item $\{0,1,\ldots,V-1\}^*$: any-length sequence of valid token IDs.
\end{itemize}

In plain words: tokenizer is a translator between human text and numbers the model can process.

\subsection{Why subword tokenization}
Subwords avoid the worst of both extremes:
\begin{itemize}[leftmargin=1.5em]
    \item Word-level tokenization: huge OOV (Out-Of-Vocabulary) risk
    \item Character-level tokenization: very long sequences and weak semantic chunking
\end{itemize}

\subsection{Effect on model training}
Tokenizer quality impacts sequence length $T$, vocabulary size $V$, and frequency distribution.
Since attention complexity is roughly $\mathcal{O}(T^2)$, better tokenization can reduce training cost and stabilize gradients.

\section{Transformer Model Theory and Equations}
\subsection{Input Representation}
For tokenized input $x \in \mathbb{N}^{B \times T}$ (B is batch size i.e. the number of independent text sequences processed simultaneously during a single forward or backward pass and T is sequence length i.e. the total number of tokens contained in each of those sequences):
\begin{itemize}[leftmargin=1.5em]
    \item token embedding: $E_{tok}(x)$
    \item position embedding: $E_{pos}(0,\ldots,T-1)$
\end{itemize}
Combined representation:
\[
H^{(0)} = E_{tok}(x) + E_{pos}
\]

\subsubsection{Terms in $H^{(0)} = E_{tok}(x) + E_{pos}$}
\begin{itemize}[leftmargin=1.5em]
    \item $x$: input token IDs.
    \item $E_{tok}(x)$: meaning vectors for those tokens.
    \item $E_{pos}$: position vectors (1st token, 2nd token, etc.).
    \item $H^{(0)}$: final starting representation before attention layers.
\end{itemize}

Why add them? A token like ``cell'' should mean something, but position also matters.
``Cell in battery'' and ``battery in cell'' are not the same pattern.

Derivation intuition:
\begin{enumerate}[leftmargin=1.5em]
    \item We need semantic identity (token embedding).
    \item We need order information (position embedding).
    \item Addition is a simple way to combine both in the same vector space.
\end{enumerate}

\subsection{Causal Self-Attention}
Attention scores before masking:
\[
A = \frac{QK^\top}{\sqrt{d_k}}
\]
Causal mask $M$ sets future positions to $-\infty$:
\[
\tilde{A}_{ij} =
\begin{cases}
A_{ij}, & j \le i \\
-\infty, & j > i
\end{cases}
\]
Then:
\[
\text{Attn}(Q,K,V)=\text{softmax}(\tilde{A})V
\]

\subsubsection{Detailed explanation of attention equations}
\textbf{1) Score matrix:}
\[
A = \frac{QK^\top}{\sqrt{d_k}}
\]
\begin{itemize}[leftmargin=1.5em]
    \item $Q$ (queries): what each token is ``looking for''.
    \item $K$ (keys): what each token ``offers''.
    \item $QK^\top$: similarity between every query and every key.
    \item $d_k$: key dimension.
    \item divide by $\sqrt{d_k}$: keeps values numerically stable so softmax does not saturate too hard.
\end{itemize}

\textbf{2) Causal masking:}
Future positions are set to $-\infty$, which means after softmax they get probability near zero.
This enforces ``no peeking ahead'' during language modeling.

\textbf{3) Weighted value mix:}
\[
\text{Attn}(Q,K,V)=\text{softmax}(\tilde{A})V
\]
\begin{itemize}[leftmargin=1.5em]
    \item softmax turns each row of scores into probabilities.
    \item $V$ (values) are the information vectors being blended.
    \item output is a weighted average of value vectors.
\end{itemize}

Derivation logic in words: compare tokens, normalize comparisons into weights, then combine information according to those weights.

\subsection{Transformer Block (conceptual)}
Each encoder layer internally follows the standard pattern:
\begin{align*}
U &= \text{LayerNorm}(H + \text{SelfAttn}(H)) \\
H' &= \text{LayerNorm}(U + \text{FFN}(U))
\end{align*}
where FFN (feed-forward network) is a position-wise non-linear projection stack.

\subsubsection{Why residual and layer normalization are used}
\begin{itemize}[leftmargin=1.5em]
    \item Residual addition ($H + \text{SelfAttn}(H)$) preserves original signal and helps gradient flow.
    \item LayerNorm stabilizes scale/variance so training does not explode or vanish easily.
    \item FFN adds non-linear transformation capacity beyond attention mixing.
\end{itemize}

Layman view: attention decides ``who should listen to whom'', while FFN helps each token ``think locally'' after hearing others.

\subsection{Output Projection}
Final token logits:
\[
Z = H'W_o + b_o, \quad Z \in \mathbb{R}^{B \times T \times V}
\]

\subsubsection{Symbols in output projection}
\begin{itemize}[leftmargin=1.5em]
    \item $H'$: final hidden features from Transformer layers.
    \item $W_o$: matrix that maps hidden features to vocabulary scores.
    \item $b_o$: bias vector.
    \item $Z$: logits (unnormalized scores) for every possible next token.
    \item $B$: batch size, $T$: sequence length, $V$: vocabulary size.
\end{itemize}

Purpose: convert internal representation into a score for each wordpiece token.

\subsection{Sequence-length Guard}
The forward pass enforces:
\[
T \leq T_{max}
\]
If violated, the implementation raises an exception, preventing out-of-range positional indexing.

\section{Training Objective and Optimization Theory}
\subsection{Language Modeling Target Shift}
Given token stream $(x_1, x_2, \ldots, x_T)$:
\begin{itemize}[leftmargin=1.5em]
    \item model input at step $t$: $x_{\le t}$
    \item target at step $t$: $x_{t+1}$
\end{itemize}

\subsection{Likelihood Objective}
\[
\max_\theta \sum_{t=1}^{T-1} \log p_\theta(x_{t+1}\mid x_{\le t})
\]
Equivalent minimization of cross-entropy loss:
\[
\mathcal{L}(\theta) = -\sum_{t=1}^{T-1}\log p_\theta(x_{t+1}\mid x_{\le t})
\]

\subsubsection{Derivation of the loss function}
At each step, the model predicts a probability for the true next token.
If the model is confident and correct, probability is high and penalty is low.
If it is wrong/confused, probability is low and penalty is high.

Why negative log?
\begin{itemize}[leftmargin=1.5em]
    \item probabilities multiply across sequence positions,
    \item log turns multiplication into addition,
    \item negative sign turns ``maximize probability'' into ``minimize loss''.
\end{itemize}

Term-by-term:
\begin{itemize}[leftmargin=1.5em]
    \item $\theta$: all model parameters.
    \item $x_{\le t}$: context up to time $t$.
    \item $x_{t+1}$: the true next token.
    \item $p_\theta(\cdot)$: predicted probability under current model.
\end{itemize}

\subsection{Batch and Tensor Shapes}
For each training segment:
\begin{itemize}[leftmargin=1.5em]
    \item Input shape: $[B,T]$
    \item Logits shape: $[B,T,V]$
    \item Flattened for CE: $[(B\cdot T),V]$ vs targets $[(B\cdot T)]$
\end{itemize}

\subsection{Optimizer}
Adam updates are applied per parameter $w$ using moment estimates:
\begin{align*}
m_t &= \beta_1 m_{t-1} + (1-\beta_1)g_t \\
v_t &= \beta_2 v_{t-1} + (1-\beta_2)g_t^2 \\
\hat{m}_t &= \frac{m_t}{1-\beta_1^t},\quad \hat{v}_t=\frac{v_t}{1-\beta_2^t} \\
w_t &= w_{t-1} - \eta\frac{\hat{m}_t}{\sqrt{\hat{v}_t}+\epsilon}
\end{align*}

\subsubsection{Adam equations explained step-by-step}
\begin{itemize}[leftmargin=1.5em]
    \item $g_t$: current gradient (direction of error reduction).
    \item $m_t$: moving average of gradients (momentum-like direction memory).
    \item $v_t$: moving average of squared gradients (scale memory).
    \item $\hat{m}_t, \hat{v}_t$: bias-corrected estimates (important at early steps).
    \item $\eta$: learning rate.
    \item $\epsilon$: small constant to avoid divide-by-zero.
\end{itemize}

Intuition:
\begin{enumerate}[leftmargin=1.5em]
    \item Keep track of trend direction ($m_t$).
    \item Keep track of how noisy/large gradients are ($v_t$).
    \item Take bigger steps where safe, smaller steps where unstable.
\end{enumerate}

\subsection{EOS (End of Sequence) Delimiting Strategy}
EOS token insertion between records introduces explicit sample boundaries in token stream, helping model avoid leakage from one QA pair into another.

\section{Inference: Generation, Scoring, and Fallback}
\subsection{Prompt Construction}
Inference prompt is consistently formatted as:
\begin{verbatim}
Ask: <user question>
Answer:
\end{verbatim}
This mirrors training format and reduces prompt-distribution shift.

\subsection{Autoregressive Decoding}
At decode step $k$:
\begin{enumerate}[leftmargin=1.5em]
    \item Build context tokens from prompt plus generated suffix.
    \item Truncate to model max context if necessary.
    \item Compute logits for last position.
    \item Apply repetition penalty to recently used tokens.
    \item Apply temperature scaling:
    \[
    z'_i = \frac{z_i}{T_{temp}}
    \]
    \item Apply top-$k$ filtering and sample.
\end{enumerate}

\subsubsection{Why temperature and top-$k$ are needed}
\begin{itemize}[leftmargin=1.5em]
    \item \textbf{Temperature}: controls randomness. Lower values make output more deterministic.
    \item \textbf{Top-$k$}: restricts sampling to best $k$ candidates, reducing absurd low-probability picks.
\end{itemize}

Analogy: temperature changes how adventurous the model is; top-$k$ limits choices to a shortlist.

\subsection{Weak-Answer Heuristics}
Generated output is marked weak if conditions such as these are met:
\begin{itemize}[leftmargin=1.5em]
    \item too short,
    \item malformed meta-markers,
    \item poor lexical overlap with question keywords.
\end{itemize}

\subsection{Retrieval Scoring Theory}
For candidate sentence $s$ and question $q$:
\begin{align*}
\text{score}(q,s) &= \text{overlap}(q,s) + \text{bonus}(q,s)
\end{align*}
where overlap is keyword match count and bonus includes definition-style and density adjustments.

\subsubsection{Scoring terms in plain language}
\begin{itemize}[leftmargin=1.5em]
    \item overlap: how many important question words appear in the candidate sentence.
    \item bonus: extra points for sentence forms that look like direct definitions/explanations.
\end{itemize}

Purpose: choose sentences that are both relevant and concise when generation is weak.

\subsection{Decision Rule}
Let $G$ be generated answer and $R$ be top retrieved sentence:
\begin{itemize}[leftmargin=1.5em]
    \item If the question is factual (for example: who/where/when/which/how many), prefer retrieval directly.
    \item If $G$ is weak, choose $R$.
    \item Else if $\text{score}(R)$ exceeds $\text{score}(G)$ by margin, choose $R$.
    \item Else for definition/open-ended questions, prefer retrieval when its score is at least comparable.
    \item Otherwise keep $G$.
\end{itemize}

\subsubsection{Why a hybrid decision is better than generation-only}
With small models, generation can be fluent but wrong, or correct but noisy.
Retrieval-only can be accurate but rigid.
The decision rule combines both:
\begin{itemize}[leftmargin=1.5em]
    \item trust generation when it is coherent and relevant,
    \item otherwise switch to retrieval for factual stability.
\end{itemize}

\subsection{Output Normalization}
Final text is cleaned for OCR artifacts, repeated spacing, and noisy prefixes to improve readability.

\section{Complexity and Resource Considerations}
\subsection{Training Complexity}
Dominant Transformer attention term is approximately:
\[
\mathcal{O}(B\cdot T^2 \cdot d)
\]
with additional feed-forward cost roughly:
\[
\mathcal{O}(B\cdot T\cdot d\cdot d_{ff})
\]

\subsection{Inference Complexity}
Autoregressive decoding over $K$ generated tokens requires repeated forward passes, making total compute roughly proportional to:
\[
\sum_{k=1}^{K} \mathcal{O}((T_0+k)^2)
\]
under full-sequence recomputation.

\subsection{Interpretation of complexity}
The $T^2$ term means attention cost grows quickly with sequence length.
If you double sequence length, attention work is roughly four times larger.
That is why good tokenization and sensible context length are crucial for local hardware.

\section{Robustness Engineering Details}
\subsection{Environment Consistency}
Pipeline scripts prefer project-local virtual environment executables when available, reducing cross-machine dependency drift.

\subsection{Fallback Layers}
Multiple fallback layers are intentionally used:
\begin{itemize}[leftmargin=1.5em]
    \item Sentence splitting fallback if NLTK resources fail (without online downloads)
    \item Retrieval fallback if generation quality is low
    \item Text cleanup fallback for OCR artifacts
\end{itemize}

\section{Step-by-Step Operational Guide}
\subsection{Full Pipeline Command}
\begin{verbatim}
.\.venv\Scripts\python.exe run_pipeline.py --epochs 30
\end{verbatim}

Optional fast-refresh run (skip tokenizer + training):
\begin{verbatim}
.\.venv\Scripts\python.exe run_pipeline.py --skip-tokenizer --skip-train
\end{verbatim}

\subsection{Manual Pipeline Commands}
\begin{verbatim}
.\.venv\Scripts\python.exe textextraction.py
.\.venv\Scripts\python.exe clean_text.py
.\.venv\Scripts\python.exe sentence_split.py
.\.venv\Scripts\python.exe create_dataset.py
.\.venv\Scripts\python.exe train_tokenizer.py
$env:EPOCHS='30'; .\.venv\Scripts\python.exe train.py
.\.venv\Scripts\python.exe generate.py
\end{verbatim}

\subsection{Iteration Cycle}
Whenever new PDFs are added:
\begin{enumerate}[leftmargin=1.5em]
    \item Re-run extraction through dataset creation
    \item Re-train tokenizer (recommended when vocabulary shifts)
    \item Re-train model checkpoint
    \item Re-evaluate sample questions
\end{enumerate}

\section{Error Analysis and Typical Failure Modes}
\subsection{Observed Failure Types}
\begin{itemize}[leftmargin=1.5em]
    \item Heading leakage into answers
    \item Punctuation corruption
    \item Over-generic generated text
    \item Prompt-template leakage (``Ask:'', ``Answer:'')
\end{itemize}

\subsection{Mitigation in Current System}
\begin{itemize}[leftmargin=1.5em]
    \item retrieval rank scoring and fallback
    \item weak-answer filtering
    \item normalization and trimming functions
    \item structured prompt consistency across train and infer
\end{itemize}

\section{Validation Strategy}
Recommended validation should include:
\begin{enumerate}[leftmargin=1.5em]
    \item fixed benchmark question set,
    \item qualitative factuality checks,
    \item comparison of generated vs retrieved outputs,
    \item regression checks after each data/model update.
\end{enumerate}

\subsection{Practical evaluation template for this project}
For each model version, record:
\begin{itemize}[leftmargin=1.5em]
    \item question,
    \item generated answer,
    \item retrieved fallback answer,
    \item final selected answer,
    \item human rating (correct / partially correct / incorrect).
\end{itemize}

This creates a simple but rigorous feedback loop for iterative improvement.

\section{Future Technical Extensions}
\begin{enumerate}[leftmargin=1.5em]
    \item Replace sentence-level retrieval with chunk-level retrieval and reranking.
    \item Add mixed-precision and gradient accumulation for larger batch simulation.
    \item Developing a graphical Web UI to replace the CLI.
\end{enumerate}

\section{Conclusion}
This project is a complete local Transformer base Small Language Model with a practical hybrid inference design.
Its technical strength lies in explicit pipeline stages, deterministic scripting, and retrieval-augmented reliability under small-model constraints.
With additional data quality and evaluation tooling, it can evolve into a stronger QA platform while remaining fully local and controllable.

\end{document}
